\scanpage{214}

\mspage{88}

(7.1.4) $\left\{\vphantom{\dfrac{A}{A}}\right.$ a) $W$ est $L$-accessible (i.e. stable par $\varinjlim$ filtr., et accessible)

b) Soient $\alpha,\beta$ composables, avec $\unsure{\beta\alpha|\beta}$ et $\alpha\in W$, alors $\unsure{\beta|\alpha\beta|\beta\alpha}\in W$. (\uline{NB} Cette condition sera satisfaite si $W$ est un localiseur modérément saturé, et c'est à cette situation que nous allons appliquer notre énoncé).

c) $\mathrm{Mon}\cap W$ est stable par cochangement de base,

% margin note written sideways
\inserted{et que $M$ satisfait les conditions supplémentaires (7.1.5) $\left\{\vphantom{\dfrac{A}{A}}\right.$ e$'$) \unsure{les} \uline{$\lim$} filtrantes \unsure{commutent} au chgt. de base f) la condition d'exactitude (7.2.1) ci-dessous}

Soit $\Pi$ un cardinal tel que $W$ soit $\Pi$-accessible, et $M_0\subset\mathrm{Filt}^{\Pi}$. (donc $M_0$ essentiellement petite). Alors on a

\[
W\cap\mathrm{Mon}\subset\widetilde{(W\cap\mathrm{Mon})^\Pi}\tag*{(7.1.6)}
\]

(ici pour \deleted{tout} partie $C$ de $\mathrm{Fl}\,M$, on pose $C^\Pi=C\cap\mathrm{Fl}^{\Pi}M$), donc

\[
\widetilde{W\cap\mathrm{Mon}}=\widetilde{(W\cap\mathrm{Mon})^\Pi},\tag*{(7.1.7)}
\]

donc $\widetilde{W\cap\mathrm{Mon}}$ est $Q$-\uline{engendré} par une \uline{petite} partie de $W\cap\mathrm{Mon}$.

\uline{Corollaire} 7.2 : Soit $M$ satisfaisant a) b) c) d) e), et \unsure{aux} conditions supplémentaires \deleted{que si}

(7.2.1) $\left\{\vphantom{\dfrac{A}{A}}\right.$ \inserted{e$'$) les \uline{$\lim$} filtr. \deleted{\unsure{st}} commutent au chgt. de base} f) Si $A_1,A_2$ sont deux sous-objets d'un objet $B$ de $M$, \deleted{alors} $A_0=A_1\cap A_2$, alors $A_1\amalg_{A_0} A_2\hookrightarrow B$ est un monomorphisme

% slanted
\margin{Une autre façon de le dire : Si $C=\mathop{\mathrm{Sup}}\limits_{\mathrm{s.o.}(B)}(A_1,A_2)$, \deleted{\unsure{alors}} \unsure{ce carré cart.|carré cart.} \diagram{\begin{tikzcd}[ampersand replacement=\&,column sep=tiny,row sep=tiny]
A_0 \arrow[r,dash] \arrow[d] \& A_1 \arrow[d,dash]\\
A_2 \arrow[r,dash] \& C
\end{tikzcd}} est aussi cocart.}

\deleted{(7.2.2) g) $\mathrm{Mon}$ est stable par \unsure{cochgt. de base}.} \deleted{Soit} \unsure{$W$} une partie de $\mathrm{Fl}\,M$, satisfaisant \deleted{localiseur mod. saturé \unsure{dans}} \deleted{$M$, tel que} conditions suivantes

(7.2.3) $\left\{\vphantom{\dfrac{A}{A}}\right.$ a) $W$ est un localiseur modérément saturé \deleted{\unsure{$M$}} b) $W$ est $L$-accessible, i.e. \inserted{Il est} stable par \uline{$\lim$} filtrantes, et est accessible

c) $\mathrm{Mon}\cap W$ \deleted{\unsure{stable}}, \inserted{stable par cochgt de base, i.e. si on a un} diagramme \inserted{\uline{cocartésien}}

\diagram{\begin{tikzcd}[ampersand replacement=\&,row sep=small]
 A \arrow[r,hook,"i"] \arrow[d,"\alpha"'] \arrow[dr,phantom,"\mathrm{coc.}"] \& B \arrow[d,"\beta"]\\
 A' \arrow[r,"i'"'] \& B'
\end{tikzcd}}

avec $i$ mono\inserted{ et $\in W$}, alors \deleted{$\alpha\in W$ implique $\beta\in W$.} \inserted{$i'$ est mono et $\in W$.}

% slanted
\margin{Cf plus bas pour une quatrième condition d) sur $W$, qui assure que la situation donne un triple de Quillen clos}
% archive sheet number: 213
